\section*{第 \ref{ch:g7:negatives} 章 --- 负数}

\begin{solution}{exo:g7:negatives:1}
相反数：$-7$；\ $+3.5$；\ $0$（它自身的相反数）；\ $-12$。两个
数互为相反数，当且仅当它们的和为 $0$。
\end{solution}

\begin{solution}{exo:g7:negatives:2}
$-4 < 1$; \quad $-4 > -7$; \quad $-2.5 < -2.4$; \quad $0 > -0.1$.
\end{solution}

\begin{solution}{exo:g7:negatives:3}
$-3.1 < -3 < -0.5 < 0 < 1 < 2.5$.
\end{solution}

\begin{solution}{exo:g7:negatives:4}
$(-5) + (-9) = -14$; \quad
$(-12) + (+7) = -5$; \quad
$(+3.5) + (-1.5) = +2$; \quad
$(-6) + (+6) = 0$.
\end{solution}

\begin{solution}{exo:g7:negatives:5}
$4 - 11 = 4 + (-11) = -7$.

$(-5) - (+3) = (-5) + (-3) = -8$.

$(-7) - (-10) = (-7) + (+10) = +3$.

$0 - (-8) = 0 + (+8) = +8$.
\end{solution}

\begin{solution}{exo:g7:negatives:6}
$(+7) + (-12) + (+3) + (-8) = (+10) + (-20) = -10$（正数
相加，负数相加）。

$(-5.5) + (+9) + (-4.5) + (+1) = (-10) + (+10) = 0$.
\end{solution}

\begin{solution}{exo:g7:negatives:7}
$(-3) + 8 - 12 = 5 - 12 = -7$：夜间温度为 $-7$ 度。
\end{solution}

\begin{solution}{exo:g7:negatives:8}
任意作图；点 $E(0, -3)$ 位于纵轴上（其横坐标
为 $0$）。
\end{solution}

\begin{solution}{exo:g7:negatives:9}
莫斯科：$(-4) - (-13) = -4 + 13 = 9$，温差为 9 度。
赫尔辛基：$2 - (-9) = 11$，温差为 11 度。
\end{solution}

\begin{solution}{exo:g7:negatives:10}
$(-12) + 7 - 4 = -5 - 4 = -9$：她位于 $-9$ 米处。要到达水面
（海拔为 $0$），她需上升 $0 - (-9) = 9$ 米。
\end{solution}

\begin{solution}{exo:g7:negatives:11}
$D(-2, -2)$：横坐标与 $A$ 相同，纵坐标与
$C$ 相同。边长：$AB = 4$（从 $-2$ 到 $2$），$BC = 3$（从 $1$ 向下到
$-2$），所以周长为 $2 \times (4 + 3) = 14$。
\end{solution}

\begin{solution}{exo:g7:negatives:12}
完整的对角线 $0, -2, -4$ 给出幻和
$0 + (-2) + (-4) = -6$。然后逐格计算：
左上格 $= -6 - (-7) - 0 = 1$；中间行右格 $= -6 - (-3) - (-2) =
-1$；右下格 $= -6 - 0 - (-1) = -5$；下中格
$= -6 - (-7) - (-2) = 3$。
\begin{center}
\renewcommand{\arraystretch}{1.4}
\begin{tabular}{|c|c|c|}
\hline
$1$ & $-7$ & $0$ \\
\hline
$-3$ & $-2$ & $-1$ \\
\hline
$-4$ & $3$ & $-5$ \\
\hline
\end{tabular}
\end{center}
每行、每列及每条对角线上的数之和都为 $-6$。
\end{solution}

\begin{solution}{pb:g7:negatives:1}
\textbf{1.} $4\,808 - (-430) = 4\,808 + 430 = 5\,238$ 米
（\cref{thm:g7:negatives:sub}：减去一个负数等于
加上它的相反数）。

\textbf{2.} 结余：
\[
(-45) + 120 + (-60) + 15 = 75 + (-60) + 15 = 15 + 15 = 30
\text{ 欧元}.
\]
要将最初的 $-45$ 恢复为 $0$，应存入其相反数：
$45$ 欧元。

\textbf{3.} $(-44) - (-100) = (-44) + 100 = 56$：凯撒生活了
$56$ 年。（这里没有陷阱：两个年份都位于
纪年分界的同一侧。）

\textbf{4.} 年份 $0$：从 $-63$ 数到 $+14$ 经过
$0$，但史学家的历法直接从
公元前 $1$ 年跳到公元 $1$ 年 --- 年份 $0$ 从未存在，因此朴素的
计数多算了一年。

\textbf{5.} 公元前 $63$ 年记为 $-(63 - 1) = -62$。然后
$14 - (-62) = 14 + 62 = 76$ 年，恰好是史学家给出的
数字。在天文学家重新编号的数轴上不存在空缺，
因此普通减法即可得到真实的时长。

\textbf{6.} 公元前 $1$ 年 $\to 0$； \ 公元前 $10$ 年 $\to -9$； \ 公元
$476 \to +476$； \ 公元前 $753$ 年 $\to -752$。

\textbf{7.} $476 - (-752) = 476 + 752 = 1\,228$ 年。

\textbf{8.} 公元前 $12$ 年即 $-11$。下一次出现的年份：
$-11 + 76 = 65$，其后为 $141$，再后为 $217$：即公元 $65$ 年、公元
$141$ 年和公元 $217$ 年。往前推：$-11 - 76 = -87$，即
历史学家所说的公元前 $87 + 1 = 88$ 年。

\textbf{9.} 公元前 $5$ 年即 $-4$，因此 $5 - (-4) = 9$ 年 --- 而非
$10$。缺失的零年每当时长跨越
纪年分界时都会产生影响。

\textbf{10.} 公元前 $3$ 年即 $-2$；十年后为 $-2 + 10 = +8$：
这个孩子在公元 $8$ 年满十岁。

\textbf{11.} 在 $-7$ 与 $-2$ 之间：$(-2) - (-7) = 5$。在
$-3.5$ 与 $4.5$ 之间：$4.5 - (-3.5) = 8$。

\textbf{12.} 将相反数及便于相加的数分组：
\[
3 + (-5) + 2 + (-1) + 4
= \bigl(3 + 2\bigr) + (-5) + (-1) + 4
= 0 + (-1) + 4 = +3 :
\]
电梯最终停在 $3$ 层。顺序并不影响结果：
正负数之和可以任意重新组合（每一步都沿同一条数轴行走，
总路程不变）。

\textbf{13.} 净变化：
$2 + (-4) + 1 + 0 + (-3) + 5 + (-2) = -1$ 度。从
$-2$ 度出发：最终温度为 $-2 + (-1) = -3$ 度。一般地，
末值等于初值加上所有的
变化量之和 --- 无需逐次跟踪升降。

\textbf{14.} 回到起点意味着净变化为 $0$：升降
必须恰好抵消。已知的五次变化之和为
$4 + (-7) + 2 + 3 + (-5) = -3$，因此第六次变化量是
其相反数：$+3$（上升三百米）。

\textbf{15.} 天文学家：公元前 $509$ 年 $\to -508$，而
$476 - (-508) = 984$ 年。史学家的 $509 + 476 = 985$
多算了一年 --- 又是那个虚构的零年，
在加法中悄然计入，却不存在于历史之中。
（作为检验，两个时期：共和国时期从 $-508$ 到 $-26$，共 $482$
年；帝国时期从 $-26$ 到 $+476$，共 $502$ 年；
$482 + 502 = 984$。）
\end{solution}
